\chapter*{Resumen}
Este proyecto desarrolla un ecosistema digital para la monitorización glucémica y la educación terapéutica mediante la integración de tecnología NFC (Near Field Communication) y Sistemas de Soporte a la Decisión (CDSS). El sistema permite la adquisición de datos de glucosa intersticial en tiempo real al conectarse directamente con sensores de monitorización continua (FreeStyle Libre 2 Plus), centralizando registros críticos como el perfil insulínico, las pautas médicas, la ingesta de carbohidratos y la actividad física.

El núcleo innovador del proyecto reside en una interfaz conversacional basada en Inteligencia Artificial Generativa (modelo tipo GPT). Este módulo actúa como un asistente pedagógico que analiza variables en tiempo real, como la Insulina Activa (IOB), para ofrecer orientaciones personalizadas ante situaciones de riesgo, incluyendo la prevención de hipoglucemias tras el ejercicio o la gestión de la insulina basal.

Bajo un enfoque estrictamente educativo, el objetivo principal es empoderar al usuario sin conocimientos previos, permitiéndole afianzar conceptos y aprender de forma intuitiva a través de sus propios datos. El sistema prioriza la seguridad, contextualizando cada respuesta de la IA como un soporte complementario que refuerza, y no sustituye, las pautas médicas establecidas. En definitiva, el proyecto demuestra cómo la convergencia de la monitorización biométrica y la IA facilita una comprensión más profunda y accesible de la gestión diaria de la diabetes.
\vspace{.5cm}

\textbf{Palabras clave:} Diabetes Mellitus, Inteligencia Artificial Generativa, Monitorización Continua de Glucosa (MCG), NFC, Sistemas de Soporte a la Decisión Clínica, Educación Terapéutica, mHealth.